(1) 无损、各向同性、互易光学活性介质在沿正 \(z\) 传播时具有右/左圆偏振本征波数 \(k_\pm(\omega)=\frac{\omega}{c}[n(\omega)\pm\kappa(\omega)]\)，其中实数 \(\kappa\) 表示弱手性折射率分裂。长度为 \(L\) 的样品输入 \(x\) 方向线偏振单色光。按明确的圆偏振基约定，推导输出Jones矢量及旋光角。
(2) 求输出Stokes参数 \(S_0,S_1,S_2,S_3\) 与旋角之间的关系，证明无吸收时单频输出仍为完全线偏振；解释若改变 \(\kappa\) 符号旋向如何变化。
(3) 入射脉冲两个圆偏振分量包络为同宽高斯 \(f(t)=\exp[-t^2/(4\sigma_t^2)]\)，且样品内频散在中心频率附近可线性展开。求两分量群延迟差 \(\Delta\tau\) 及时间积分后的交叉相干可见度因子 \(|\mathcal V|\)，推导何时即使无吸收也会观察到曝光平均退偏。
(4) 样品两端配一面普通镜，将光往返传播一次。利用互易介质反向圆偏振本征标记的交换，解释在理想一致路径与无镜面额外各向异性时天然旋光在往返中抵消；与非互易Faraday旋转通常累积的情况对照，并检验 \(\kappa=0\) 和 \(\Delta\tau=0\) 极限。
